\documentclass[11pt]{article}
\usepackage[margin=1.05in]{geometry}
\usepackage{amsmath,amssymb,amsthm,mathtools}
\usepackage{hyperref}

\newtheorem{theorem}{Theorem}[section]
\newtheorem{proposition}[theorem]{Proposition}
\newtheorem{lemma}[theorem]{Lemma}
\newtheorem{corollary}[theorem]{Corollary}
\theoremstyle{definition}
\newtheorem{remark}[theorem]{Remark}
\newtheorem{conjecture}[theorem]{Conjecture}

\newcommand{\Z}{\mathbb Z}
\newcommand{\R}{\mathbb R}
\newcommand{\PFtwo}{\mathrm{PF}_2}
\newcommand{\Bx}{\mathrm{Box}}
\newcommand{\Ylat}{\mathcal Y}
\DeclarePairedDelimiter{\abs}{\lvert}{\rvert}

\title{The inter-region inequality at $m=2$ is false,\\
and condition (A) at $m=2$ is true}
\author{Clio}
\date{1 October 2026}

\begin{document}
\maketitle

\begin{abstract}
The brief for this session asked for a proof of the inter-region inequality~(G), the
single remaining gap in \texttt{2026-09-30-c1-cylindric-kostka-logconcavity.tex}. Step~0 of
the brief --- extend the verification range before proving anything --- refutes it.
\textbf{(G) is false}: at $n=10$, $\mu=(0,5)$, $\lambda=(5,12)$, $b=4$ the region
partial sums $G_{\mathrm I}=(1,3,6,7,6,3,1)$ and $G_{\mathrm{III}}=(1,1,1,1,1,1,1)$ share the
support $[10,16]$ and the pairwise inequality reads $6\ge 7$ at $s=11$. The failure is
structural, not marginal: against a partner that is \emph{constant} across one's support,
the pairwise inequality~$(*)$ is \emph{literally concavity}, and concavity of the region
sums was already proved false four pages later in the same paper. The route was refuted by
a fact the paper itself contains.

Chasing the mechanism, however, produced the theorem. Two one-line identities,
\[
  a_t+d_t=t+D,\qquad b_t+c_t=u+B-t,
\]
force \emph{every} trapezoid in the slice sum to be symmetric about the \emph{same} centre
$\sigma_0=\tfrac12(u+\lambda_1+\lambda_2)$. Hence $G$ is a radial function,
$G(s)=\gamma\bigl(\abs{s-\sigma_0}\bigr)$, where $\gamma$ is the upper tail of a sequence
$\beta$ that is the positive part of a concave function. Since the upper tail of any
$\PFtwo$ sequence is log-concave, $G$ is log-concave. This proves
\textbf{condition~(A) at $m=2$ outright} --- all slices, not $79.9\%$ --- without~(G),
and it explains in one identity why (G) and concavity both fail:
$2\gamma(y)-\gamma(y-1)-\gamma(y+1)=\beta(y)-\beta(y-1)$.
\end{abstract}

\section{What is proved and what is refuted}

Throughout, $m=2$, $\mu=(\mu_1,\mu_2)$, $\lambda=(\lambda_1,\lambda_2)$, $n\ge2$ is the
level, $d=|\lambda/\mu|$, $\ell=3$, and $k(a,b)=K^c_{(a,b,d-a-b)}$ as in
\texttt{2026-09-30-c1-cylindric-kostka-logconcavity.tex} (henceforth \textbf{[c1]}). Write
$\PFtwo$ for: nonnegative, support an interval of $\Z$, and $g(s)^2\ge g(s-1)g(s+1)$ for
all $s$.

\begin{theorem}[Condition (A) at $m=2$]\label{thm:main}
For every cylindric skew shape $\lambda/\mu$ with $m=2$, every level $n\ge2$ and every $b$,
the sequence $a\mapsto k(a,b)$ is $\PFtwo$. In particular
$k(a,b)^2\ge k(a+1,b)k(a-1,b)$ for all $a,b$: condition~(A) of
\textup{[c1, Conj.~2.1]} holds at $m=2$.
\end{theorem}

\begin{proposition}[the inter-region inequality is false]\label{prop:Gfalse}
Inequality~(G) of \textup{[c1, §"The gap"]} --- that
\[
 2G_{\mathrm R}(s)G_{\mathrm R'}(s)\ \ge\ G_{\mathrm R}(s-1)G_{\mathrm R'}(s+1)
   +G_{\mathrm R}(s+1)G_{\mathrm R'}(s-1)\qquad\text{for all }s \tag{G}
\]
for every pair of distinct regions --- is false. Over $2\le n\le 11$, $m=2$, $d\le14$
($201\,147$ nonempty slices, $37\,672$ occurring region pairs) it fails on exactly $92$
pairs; every failure is the pair $\mathrm I$--$\mathrm{III}$; and every failure has all
three of $\mathrm I$, $\mathrm{II}/\mathrm{IV}$, $\mathrm{III}$ nonempty.
\end{proposition}

Theorem~\ref{thm:main} does not go through (G): it replaces the whole region decomposition.
Section~\ref{sec:refute} records the refutation and its mechanism; §\ref{sec:conc}--%
§\ref{sec:main} prove Theorem~\ref{thm:main}; §\ref{sec:why} shows that the mechanism of
the refutation is an \emph{identity} in the new picture, so (G) is dead with no residual
scope.

\section{The refutation}\label{sec:refute}

\subsection{Setting, recalled so that nothing is reconstructed from memory}

Put $A=\lambda_2-n+1$, $B=\lambda_1$, $C=\lambda_1+1$, $D=\lambda_2$, so that
\begin{equation}\label{eq:AD}
  C=B+1,\qquad A=D-n+1 .
\end{equation}
Fix $b$ and set $u=|\mu|+b$. By [c1, Prop.~2.3] (\texttt{A-slice-sum-reformulation},
\texttt{proved}), $k(a,b)=G(u+a)$ where
\begin{equation}\label{eq:G}
  G(s)=\sum_{t\in T}\bigl(\mathbf 1_{[a_t,b_t]}*\mathbf 1_{[c_t,d_t]}\bigr)(s),
  \qquad T=[T_-,T_+],
\end{equation}
\begin{equation}\label{eq:endpoints}
  a_t=\max(t,A),\quad b_t=\min(u-1-t,B),\quad c_t=\max(u-t,C),\quad d_t=\min(t+n-1,D),
\end{equation}
and $T_-=\max(\mu_1,u-\mu_1-n+1)$, $T_+=\min(\mu_2-1,u-\mu_2)$. Here $t=\nu_1$ indexes the
first interior layer of the chain, $\nu_2=u-t$, and $s=\kappa_1+\kappa_2$. Values of $t$ for
which $a_t>b_t$ or $c_t>d_t$ contribute $0$. Write
\[
  n_t=b_t-a_t+1,\qquad m_t=d_t-c_t+1,\qquad h_t=\min(n_t,m_t).
\]
The two breakpoints are $\tau_1=A$ (for $a_t$ and $d_t$) and $\tau_2=u-1-B=u-C$ (for $b_t$
and $c_t$), and the four regions of [c1] are indexed by
$\bigl(\mathbf 1[t\ge\tau_1],\ \mathbf 1[t>\tau_2]\bigr)=(0,0),(1,0),(1,1),(0,1)$ for
$\mathrm I,\mathrm{II},\mathrm{III},\mathrm{IV}$ respectively, with $G_{\mathrm R}$ the
corresponding partial sum of~\eqref{eq:G}.

\subsection{Only three regions ever occur, and they are a chain}

\begin{lemma}\label{lem:threeregions}
$T_{\mathrm{II}}$ and $T_{\mathrm{IV}}$ are never both nonempty. Consequently at most three
regions occur, and in increasing order of $t$ they are
$\mathrm I,\ \mathrm{II},\ \mathrm{III}$ \textup{(if $\tau_1\le\tau_2$)} or
$\mathrm I,\ \mathrm{IV},\ \mathrm{III}$ \textup{(if $\tau_2<\tau_1$)}.
\end{lemma}
\begin{proof}
$T_{\mathrm{II}}=\{t\in T:\tau_1\le t\le\tau_2\}\ne\emptyset$ forces $\tau_1\le\tau_2$,
while $T_{\mathrm{IV}}=\{t\in T:\tau_2<t<\tau_1\}\ne\emptyset$ forces $\tau_2<\tau_1$.
The order statement is immediate from the defining inequalities.
\end{proof}

This already corrects the brief, which proposed attacking the pair $\mathrm{II}$--%
$\mathrm{IV}$ first on the ground that both are $\rho\circ H$ with a shared $H$. That pair
is \emph{vacuous}; consistently, the census of [c1] lists $\mathrm{II}$--$\mathrm{IV}$
nowhere and its five pair counts sum to its total ($177+177+81+45+45=525$).

\subsection{The mechanism: against a flat partner, $(*)$ is concavity}

\begin{lemma}\label{lem:flat}
Let $g,h\ge0$ and suppose $h\equiv c>0$ on an interval $J$. If $s-1,s,s+1\in J$ then
\[
  2g(s)h(s)\ \ge\ g(s-1)h(s+1)+g(s+1)h(s-1)
  \iff
  2g(s)\ \ge\ g(s-1)+g(s+1).
\]
\end{lemma}
\begin{proof}
Both sides of the left inequality equal $c$ times the corresponding side of the right one.
\end{proof}

So whenever one region sum is constant and positive across a window containing three
consecutive points of another's support, (G) for that pair \emph{is} concavity of the other
at those points. And [c1, §"Concavity is false"] already proves concavity of the slice sums
false, with witness $k(\cdot,2)=(1,3,6,7,6,3,1)$ at $n=8$, $\mu=(0,3)$, $\lambda=(4,7)$.
All that is needed to refute (G) is a shape in which a non-concave region sum coexists with
a \emph{flat} one on the same support. It exists.

\subsection{The witness}

\begin{proposition}\label{prop:witness}
Let $n=10$, $\mu=(0,5)$, $\lambda=(5,12)$, $b=4$ \textup{(}so $d=12$, $u=9$, $A=3$, $B=5$,
$C=6$, $D=12$, $\tau_1=\tau_2=3$, $T=[0,4]$\textup{)}. Then
\[
\begin{array}{c|c|c|c}
 t & [a_t,b_t]\times[c_t,d_t] & \text{region} & \text{trapezoid on its support}\\\hline
 0 & [3,5]\times[9,9]   & \mathrm I   & (1,1,1)\ \text{on }[12,14]\\
 1 & [3,5]\times[8,10]  & \mathrm I   & (1,2,3,2,1)\ \text{on }[11,15]\\
 2 & [3,5]\times[7,11]  & \mathrm I   & (1,2,3,3,3,2,1)\ \text{on }[10,16]\\
 3 & [3,5]\times[6,12]  & \mathrm{II} & (1,2,3,3,3,3,3,2,1)\ \text{on }[9,17]\\
 4 & [4,4]\times[6,12]  & \mathrm{III}& (1,1,1,1,1,1,1)\ \text{on }[10,16]
\end{array}
\]
so that, all on $[10,16]$,
\[
  G_{\mathrm I}=(1,3,6,7,6,3,1),\qquad G_{\mathrm{III}}=(1,1,1,1,1,1,1),
\]
and $G_{\mathrm{II}}=(1,2,3,3,3,3,3,2,1)$ on $[9,17]$, $G_{\mathrm{IV}}=0$. Each
$G_{\mathrm R}$ is $\PFtwo$. At $s=11$,
\[
  2G_{\mathrm I}(11)G_{\mathrm{III}}(11)=2\cdot3\cdot1=6
  \ <\ 7=1\cdot1+6\cdot1
  =G_{\mathrm I}(10)G_{\mathrm{III}}(12)+G_{\mathrm I}(12)G_{\mathrm{III}}(10),
\]
so \textup{(G)} fails for the pair $\mathrm I$--$\mathrm{III}$ \textup{(}and again at
$s=15$, by symmetry\textup{)}. Nevertheless
$G=(1,4,7,10,11,10,7,4,1)$ on $[9,17]$ is $\PFtwo$ --- indeed concave on its support.
\end{proposition}
\begin{proof}
Direct substitution in~\eqref{eq:endpoints}; the region labels are read off from
$\mathbf 1[t\ge3]$, $\mathbf 1[t>3]$. $G_{\mathrm I}$ is the sum of the first three
trapezoids, $(1,1,1)$ on $[12,14]$ plus $(1,2,3,2,1)$ on $[11,15]$ plus
$(1,2,3,3,3,2,1)$ on $[10,16]$, which is $(1,3,6,7,6,3,1)$ on $[10,16]$. The displayed
numerical inequality is $6<7$. The region $\mathrm{III}$ sum is a single trapezoid with
$n_4=1$, hence an interval indicator, hence constant on $[10,16]$; so by
Lemma~\ref{lem:flat} the failure at $s=11$ is exactly the failure of concavity of
$G_{\mathrm I}$ there, $2\cdot3<1+6$.
\end{proof}

Observe that $(1,3,6,7,6,3,1)$ is the \emph{same} sequence that refutes concavity in
[c1, §"Concavity is false"] and in the registry node \texttt{A-concavity}; it reappears here
as a region partial sum rather than as a full slice sum. The $525/525$ sweep of [c1]
($n\le8$, $d\le10$) missed the refutation because no slice in that range places a
non-concave region sum beside a flat one; the smallest witnesses have $n=10$, $d=12$.

\begin{remark}[census]\label{rem:census}
Over $2\le n\le11$, $m=2$, $d\le14$: $201\,147$ nonempty slices; $111\,009$ meet one region,
$30\,181$ meet two, $2\,497$ meet three (and $57\,460$ have $G\equiv0$); $37\,672$ occurring
region pairs, of which $\mathrm{II}$--$\mathrm{IV}$ accounts for $0$, confirming
Lemma~\ref{lem:threeregions}. (G) fails on $92$ pairs, all $\mathrm I$--$\mathrm{III}$, all
with three regions nonempty. In every one of the $92$ failures the total $G$ is still
$\PFtwo$, and condition~(A) fails on \emph{none} of the $201\,147$ slices. In all $92$ the
mechanism of Lemma~\ref{lem:flat} applies: one of the two sums is constant across the
support of the other.
\end{remark}

\begin{corollary}
The pairwise route of \textup{[c1, eq.~$(*)$]} cannot close condition~(A) at $m=2$.
Since $(*)$ also fails between individual slice members
\textup{([c1, §2.1 item 1]: $n=6$, $\mu=(0,3)$, $\lambda=(4,5)$)}, the symmetrisation
identity is unusable at \emph{both} available granularities --- members and regions.
\end{corollary}

\section{Two identities, and concentricity}\label{sec:conc}

The refutation says the region decomposition is the wrong decomposition. The following two
identities say what the right one is. They are the entire content of the proof: everything
afterwards is bookkeeping.

\begin{lemma}[crossed sums]\label{lem:crossed}
For every $t\in\Z$,
\[
  a_t+d_t=t+D,\qquad b_t+c_t=u+B-t .
\]
\end{lemma}
\begin{proof}
By~\eqref{eq:AD}, $a_t+d_t=\max(t,A)+\min(t+n-1,\,A+n-1)$. If $t\le A$ this is
$A+(t+n-1)$; if $t\ge A$ it is $t+(A+n-1)$. Both equal $t+A+n-1=t+D$.

Again by~\eqref{eq:AD}, $b_t+c_t=\min(u-1-t,B)+\max(u-t,\,B+1)$. If $u-1-t\le B$ then
$u-t\le B+1$ and the value is $(u-1-t)+(B+1)$; if $u-1-t\ge B$ then $u-t\ge B+1$ and the
value is $B+(u-t)$. Both equal $u+B-t$.
\end{proof}

The two hypotheses $A=D-n+1$ and $C=B+1$ are each used exactly once, and each is used to
make a $\max$/$\min$ pair agree \emph{on both branches} rather than merely at the
breakpoint. This is strictly stronger than the statement in [c1] that the four endpoint
functions share only two breakpoints: it says the two \emph{crossed} sums are affine in $t$
with slopes $+1$ and $-1$ globally.

\begin{proposition}[concentricity]\label{prop:concentric}
Put
\[
  \sigma_0=\frac{u+B+D}{2}=\frac{u+\lambda_1+\lambda_2}{2},
  \qquad
  t_0=\frac{u+B-D}{2}.
\]
Then for every $t\in\Z$,
\begin{align}
  (a_t+c_t)+(b_t+d_t)&=2\sigma_0, \label{eq:center}\\
  n_t-m_t&=u+B-D-2t=2(t_0-t). \label{eq:diff}
\end{align}
Consequently every trapezoid
$\mathrm{Trap}_t:=\mathbf 1_{[a_t,b_t]}*\mathbf 1_{[c_t,d_t]}$ in~\eqref{eq:G} is symmetric
about the \emph{same} point $\sigma_0$, and $G$ is symmetric about $\sigma_0$.
\end{proposition}
\begin{proof}
$(a_t+c_t)+(b_t+d_t)=(a_t+d_t)+(b_t+c_t)=(t+D)+(u+B-t)=u+B+D$ by
Lemma~\ref{lem:crossed}, giving~\eqref{eq:center}. And
$n_t-m_t=(b_t-a_t)-(d_t-c_t)=(b_t+c_t)-(a_t+d_t)=(u+B-t)-(t+D)$,
giving~\eqref{eq:diff}.

$\mathrm{Trap}_t$ is supported on $[a_t+c_t,\ b_t+d_t]$ and, being a convolution of two
interval indicators, is symmetric about the midpoint of that interval, which
by~\eqref{eq:center} is $\sigma_0$ --- independently of $t$. A sum of functions symmetric
about a common point is symmetric about it.
\end{proof}

Note that $\sigma_0$ and $t_0$ do not depend on $t$, and that $\sigma_0-t_0=D$,
$\sigma_0+t_0=u+B$. Set
\[
  L_t:=\frac{n_t+m_t-2}{2},\qquad \delta_t:=\frac{\abs{n_t-m_t}}{2}\overset{\eqref{eq:diff}}{=}\abs{t-t_0},
\]
so that $\mathrm{Trap}_t$ is supported on $[\sigma_0-L_t,\ \sigma_0+L_t]$ and
$L_t-\delta_t=h_t-1$, $L_t+\delta_t=\max(n_t,m_t)-1$. Both $L_t$ and $\delta_t$ lie in
$\Z-t_0$; indeed $L_t=n_t+(t-t_0)-1$ by~\eqref{eq:diff}.

\begin{lemma}[the tent]\label{lem:tent}
$t\mapsto L_t$ is concave on $\Z$ with increments in $\{-1,0,1\}$; its increments are $+1$
for $t<\min(\tau_1,\tau_2)$, $0$ on $[\min(\tau_1,\tau_2),\max(\tau_1,\tau_2)]$, and $-1$
for $t\ge\max(\tau_1,\tau_2)$. Writing
$\tau_\wedge=\min(\tau_1,\tau_2)$, $\tau_\vee=\max(\tau_1,\tau_2)$ and
$\Lambda=L_{\tau_\wedge}=L_{\tau_\vee}$, we have for all $t\in\R$
\begin{equation}\label{eq:tent}
  \widetilde L(t)=\min\bigl(t-\alpha,\ \Lambda,\ \omega-t\bigr),\qquad
  \alpha=\tau_\wedge-\Lambda,\quad \omega=\tau_\vee+\Lambda,
\end{equation}
where $\widetilde L$ is the piecewise affine extension of $L$, and
\begin{equation}\label{eq:Lambda}
  \Lambda=\sigma_0-\max(u,\ A+C).
\end{equation}
\end{lemma}
\begin{proof}
$n_t=\min(u-1-t,B)-\max(t,A)+1$ is a concave function of $t$ (a concave minus a convex one)
whose increment is $-\mathbf 1[t\ge\tau_2]-\mathbf 1[t\ge\tau_1]$ (each $\min$/$\max$
contributing $-1$ past its own breakpoint); symmetrically $m_t$ is concave with increment
$\mathbf 1[t<\tau_1]+\mathbf 1[t<\tau_2]$. Hence $n_t+m_t$ is concave with increment
\[
  \bigl(\mathbf 1[t<\tau_1]-\mathbf 1[t\ge\tau_1]\bigr)
 +\bigl(\mathbf 1[t<\tau_2]-\mathbf 1[t\ge\tau_2]\bigr)\in\{-2,0,2\},
\]
equal to $+2$ iff $t<\tau_\wedge$ and to $-2$ iff $t\ge\tau_\vee$. Dividing by $2$ gives the
statement about $L$, and a concave piecewise affine function with increments $+1,0,-1$ in
that order is exactly~\eqref{eq:tent}.

For~\eqref{eq:Lambda}: by~\eqref{eq:center}, $L_t=\sigma_0-(a_t+c_t)$, so
$\Lambda=\max_tL_t=\sigma_0-\min_t(a_t+c_t)$. Now
$a_t+c_t=\max(t,A)+\max(u-t,C)$ is convex with increment $-1$ for $t<\tau_\wedge$, $0$ on
$[\tau_\wedge,\tau_\vee]$ and $+1$ for $t\ge\tau_\vee$, so its minimum is attained on
$[\tau_\wedge,\tau_\vee]$. If $\tau_1\le\tau_2$ that range is region $\mathrm{II}$, where
$a_t+c_t=t+(u-t)=u$; if $\tau_2<\tau_1$ it is region $\mathrm{IV}$, where $a_t+c_t=A+C$.
The two cases are distinguished by $\tau_1\le\tau_2\iff A+C\le u$, so in both
$\min_t(a_t+c_t)=\max(u,A+C)$.
\end{proof}

\section{The radial form}\label{sec:radial}

\begin{proposition}[radial form of the slice sum]\label{prop:radial}
Define
\begin{equation}\label{eq:beta}
  \beta(r)\ :=\ \#\{t\in T:\ \delta_t\le r\le L_t\}\qquad (r\in\Z-t_0),
\end{equation}
and $\gamma(y):=\sum_{r\ge y}\beta(r)$. Then
\begin{equation}\label{eq:radial}
  G(s)\ =\ \gamma\bigl(\abs{s-\sigma_0}\bigr)\qquad\text{for all }s\in\Z .
\end{equation}
\end{proposition}
\begin{proof}
First, $\abs{s-\sigma_0}$ and $\delta_t,L_t$ all lie in the same coset of $\Z$: indeed
$2\sigma_0=u+B+D$ and $2t_0=u+B-D$ differ by $2D$, so $\sigma_0\equiv t_0\pmod 1$, while
$\delta_t=\abs{t-t_0}\in\Z-t_0$ and $L_t=n_t+(t-t_0)-1\in\Z-t_0$. So~\eqref{eq:radial} is a
statement on one lattice.

Fix $t$ with $n_t,m_t\ge1$ and write $x=s-\sigma_0$. By Proposition~\ref{prop:concentric},
$\mathrm{Trap}_t$ is supported on $[\sigma_0-L_t,\sigma_0+L_t]$, and the convolution of
interval indicators of widths $n_t,m_t$ is
\[
  \mathrm{Trap}_t(s)=\min\bigl(s-(a_t+c_t)+1,\ (b_t+d_t)-s+1,\ n_t,\ m_t\bigr)_+
   =\min\bigl(L_t+1-\abs{x},\ h_t\bigr)_+ ,
\]
using $a_t+c_t=\sigma_0-L_t$, $b_t+d_t=\sigma_0+L_t$ and
$\min(\alpha,\beta)_+=\min(\alpha_+,\beta_+)$. Since $h_t=L_t+1-\delta_t$,
\[
  \min\bigl(L_t+1-\abs{x},\ h_t\bigr)_+
  =\bigl(L_t-\max(\abs{x},\delta_t)+1\bigr)_+
  =\#\{r:\ \max(\abs{x},\delta_t)\le r\le L_t\},
\]
the last count being over the lattice $\Z-t_0$. If instead $\min(n_t,m_t)\le0$ then
$\mathrm{Trap}_t\equiv0$, and also $h_t\le0$, i.e.\ $L_t<\delta_t$, so no $r$ satisfies
$\delta_t\le r\le L_t$ and the count is $0$ as well. Hence in all cases
\[
  \mathrm{Trap}_t(s)=\#\{r:\ \abs{x}\le r,\ \delta_t\le r\le L_t\}.
\]
Summing over $t\in T$ and exchanging the order of summation,
\[
  G(s)=\sum_{t\in T}\#\{r\ge\abs{x}:\ \delta_t\le r\le L_t\}
      =\sum_{r\ge\abs{x}}\beta(r)=\gamma(\abs{x}).\qedhere
\]
\end{proof}

\begin{proposition}[$\beta$ is $\PFtwo$]\label{prop:betaPF2}
With $\Lambda,\alpha,\omega$ as in Lemma~\ref{lem:tent} and $N=T_+-T_-+1\ge1$, put
\begin{equation}\label{eq:c}
  c(r)\ :=\ \min\Bigl(\
   \underbrace{\min(T_+,\ t_0+r,\ \omega-r)}_{\text{concave}}
  -\underbrace{\max(T_-,\ t_0-r,\ \alpha+r)}_{\text{convex}}+1,
  \ \ N(\Lambda+1-r)\ \Bigr).
\end{equation}
Then $c$ is concave and $\Z$-valued on $\Z-t_0$, and $\beta=c_+$. In particular $\beta$ is
$\PFtwo$ with finite support.
\end{proposition}
\begin{proof}
\emph{The row description.} For $t\in\Z$ we have $\delta_t\le r\iff\abs{t-t_0}\le r\iff
t_0-r\le t\le t_0+r$. By Lemma~\ref{lem:tent}, $L_t\ge r$ holds iff
$\widetilde L(t)\ge r$, which by~\eqref{eq:tent} holds iff $r\le\Lambda$ and
$\alpha+r\le t\le\omega-r$. Hence for $r\le\Lambda$
\begin{equation}\label{eq:row}
  \beta(r)=\#\bigl\{t\in\Z:\ \max(T_-,t_0-r,\alpha+r)\le t\le\min(T_+,t_0+r,\omega-r)\bigr\},
\end{equation}
and $\beta(r)=0$ for $r>\Lambda$.

\emph{Integrality.} For $r\in\Z-t_0$ we have $t_0+r\in\Z$ and $t_0-r=2t_0-(t_0+r)\in\Z$
since $2t_0=u+B-D\in\Z$. Also $\Lambda\in\Z-t_0$ (it equals $L_{\tau_\wedge}$ and
$\tau_\wedge\in\Z$), hence $\alpha=\tau_\wedge-\Lambda\in\Z+t_0$ and
$\omega=\tau_\vee+\Lambda\in\Z-t_0$, so $\alpha+r\in\Z$ and $\omega-r\in\Z$. Finally
$T_\pm\in\Z$. Thus both the $\min$ and the $\max$ in~\eqref{eq:c} are integers, and the
count in~\eqref{eq:row} equals $\bigl(\min(\cdots)-\max(\cdots)+1\bigr)_+$ with no floors.

\emph{$\beta=c_+$.} Write $c_1(r)=\min(T_+,t_0+r,\omega-r)-\max(T_-,t_0-r,\alpha+r)+1$, so
that $\beta(r)=c_1(r)_+$ for $r\le\Lambda$. Since the $\min$ includes $T_+$ and the $\max$
includes $T_-$ we have $c_1(r)\le T_+-T_-+1=N$ for every $r$. If $r\le\Lambda$ then
$\Lambda+1-r\ge1$, so $N(\Lambda+1-r)\ge N\ge c_1(r)$ and $c(r)=c_1(r)$, whence
$c(r)_+=\beta(r)$. If $r>\Lambda$ then $r\ge\Lambda+1$ (both lie in $\Z-t_0$), so
$N(\Lambda+1-r)\le0$, hence $c(r)\le0$ and $c(r)_+=0=\beta(r)$.

\emph{Concavity.} A minimum of affine functions is concave and a maximum of affine
functions is convex, so $c_1$ is concave; $N(\Lambda+1-r)$ is affine; a minimum of two
concave functions is concave. Hence $c$ is concave on $\Z-t_0$.

\emph{Conclusion.} By [c1, Lem.~2.2] (\texttt{lem:trunc}, \texttt{lean-verified} as
\verb|TworowD4Kernel.DiscreteConcavity.PFtwo_posPart|), the positive part of a concave
integer sequence has interval support and is log-concave; so $\beta$ is $\PFtwo$. Its
support is contained in $[0,\Lambda]$ --- bounded below because $\delta_t\ge0$ and above by
$\Lambda$ --- hence finite.
\end{proof}

\begin{remark}
The clamp $N(\Lambda+1-r)$ in~\eqref{eq:c} is \emph{not} cosmetic. For
$\Lambda<r\le\Lambda+\tfrac12\abs{\tau_1-\tau_2}$ the three-interval description
in~\eqref{eq:row} is non-empty while $\{t:L_t\ge r\}$ is empty, because
$\widetilde L$ has a genuine flat piece and $\min(t-\alpha,\omega-t)$ exceeds $\Lambda$
there. Dropping the clamp gives the wrong $\beta$ in $7245$ of $13\,761$ tested instances
(§\ref{sec:verif}).
\end{remark}

\section{The upper tail of a log-concave sequence}\label{sec:tail}

\begin{lemma}[tail lemma]\label{lem:tail}
Let $\beta:\Z\to[0,\infty)$ be log-concave with interval support and $\sum_r\beta(r)<\infty$.
Then $\gamma(y):=\sum_{r\ge y}\beta(r)$ satisfies
\[
  \gamma(y)^2\ \ge\ \gamma(y-1)\,\gamma(y+1)\qquad\text{for every }y\in\Z,
\]
$\gamma$ is nonincreasing, and $\{\gamma>0\}$ is an interval \textup{(}a down-set\textup{)}.
\end{lemma}
\begin{proof}
Monotonicity and the shape of $\{\gamma>0\}$ are immediate from $\beta\ge0$: $\gamma$ is
nonincreasing, and $\gamma(y)>0$ iff $y\le r_1:=\max\{r:\beta(r)>0\}$.

Fix $y$. From $\gamma(y-1)=\gamma(y)+\beta(y-1)$ and $\gamma(y+1)=\gamma(y)-\beta(y)$,
\begin{equation}\label{eq:tailid}
  \gamma(y)^2-\gamma(y-1)\gamma(y+1)
  \;=\;\beta(y-1)\beta(y)\;-\;\Delta\,\gamma(y),
  \qquad \Delta:=\beta(y-1)-\beta(y).
\end{equation}
If $\Delta\le0$ the right side of~\eqref{eq:tailid} is $\ge0$ and we are done. So assume
$\Delta>0$, i.e.\ $\beta(y-1)>\beta(y)\ge0$; we must show
$\Delta\,\gamma(y)\le\beta(y-1)\beta(y)$.

\emph{Case $\beta(y)=0$.} Then $\beta(y-1)>0$, so $y-1$ lies in the support
$[r_0,r_1]$, while $y\notin[r_0,r_1]$; since $y>y-1\ge r_0$ this forces $y>r_1$, hence
$\beta(r)=0$ for all $r\ge y$ and $\gamma(y)=0$. Both sides vanish.

\emph{Case $\beta(y)>0$.} Put $q=\beta(y)/\beta(y-1)\in(0,1)$. Log-concavity gives
$\beta(r+1)/\beta(r)\le\beta(r)/\beta(r-1)$ whenever $\beta(r)>0$; since the support is an
interval, the ratios $\beta(r+1)/\beta(r)$ for $r_0\le r<r_1$ form a nonincreasing sequence.
As $y-1\ge r_0$ (because $\beta(y-1)>0$) and the ratio at $r=y-1$ is $q$, we get
$\beta(r+1)/\beta(r)\le q$ for every $r$ with $y-1\le r<r_1$. Telescoping,
\[
  \beta(y+k)\ \le\ \beta(y)\,q^{\,k}\qquad (k\ge0),
\]
(the inequality being $0\le\beta(y)q^k$ once $y+k>r_1$). Hence
\[
  \gamma(y)=\sum_{k\ge0}\beta(y+k)\ \le\ \beta(y)\sum_{k\ge0}q^{\,k}
  =\frac{\beta(y)}{1-q}
  =\frac{\beta(y)\,\beta(y-1)}{\beta(y-1)-\beta(y)}
  =\frac{\beta(y-1)\beta(y)}{\Delta},
\]
which is exactly $\Delta\,\gamma(y)\le\beta(y-1)\beta(y)$.
\end{proof}

\begin{remark}
Lemma~\ref{lem:tail} is the discrete form of the classical fact that a log-concave
probability mass function has a log-concave survival function (equivalently, that the
log-concave class is contained in the increasing-failure-rate class). I give the proof
because it is three lines and because the hypothesis matters: \emph{interval support is
load-bearing but concavity is not}. The positive part of a concave sequence is more than
enough; mere log-concavity suffices. Dropping log-concavity while keeping interval support
destroys the conclusion in $82.6\%$ of random tests (§\ref{sec:verif}), so this is not a
hypothesis carried along for safety.
\end{remark}

\begin{lemma}[radial lemma]\label{lem:radialLC}
Let $\sigma_0\in\tfrac12\Z$ and let $\gamma$ be nonnegative and nonincreasing on
$\Ylat:=\{\abs{s-\sigma_0}:s\in\Z\}$ with $\gamma(y)^2\ge\gamma(y-1)\gamma(y+1)$ for all
$y\in\Ylat$ with $y\ge1$, and with $\{\gamma>0\}$ a down-set. Then
$g(s)=\gamma(\abs{s-\sigma_0})$ is $\PFtwo$.
\end{lemma}
\begin{proof}
Write $x=s-\sigma_0$ and note $\Ylat=\Z_{\ge0}$ if $\sigma_0\in\Z$ and
$\Ylat=\tfrac12+\Z_{\ge0}$ otherwise.
\begin{itemize}
\item If $x\ge1$ then $\abs{x-1}=x-1$ and $\abs{x+1}=x+1$, so the required inequality is
$\gamma(x)^2\ge\gamma(x-1)\gamma(x+1)$, which holds by hypothesis. If $x\le-1$ the same
argument applies to $-x$.
\item If $x=0$ (possible only when $\sigma_0\in\Z$) then $g(s\pm1)=\gamma(1)$ and
$g(s)=\gamma(0)\ge\gamma(1)\ge0$, so $g(s)^2\ge\gamma(1)^2=g(s-1)g(s+1)$.
\item If $x=\tfrac12$ (possible only when $\sigma_0\notin\Z$) then $\abs{x-1}=\tfrac12$ and
$\abs{x+1}=\tfrac32$, so the inequality reads
$\gamma(\tfrac12)^2\ge\gamma(\tfrac12)\gamma(\tfrac32)$, true since $\gamma$ is
nonincreasing and nonnegative. The case $x=-\tfrac12$ is symmetric.
\end{itemize}
Finally $\{g>0\}=\{s:\abs{s-\sigma_0}\le r_1\}$ where $r_1=\max\{\gamma>0\}$ (or
$\{g>0\}=\emptyset$), an interval of $\Z$.
\end{proof}

\section{Proof of Theorem~\ref{thm:main}}\label{sec:main}

\begin{proof}[Proof of Theorem~\ref{thm:main}]
Fix $\lambda/\mu$ with $m=2$, a level $n\ge2$ and a value of $b$. By
[c1, Prop.~2.3], $k(a,b)=G(u+a)$ with $G$ as in~\eqref{eq:G}, $u=|\mu|+b$; so it suffices to
prove $G\in\PFtwo$. If $T=\emptyset$ or every $t\in T$ has $\min(n_t,m_t)\le0$ then
$G\equiv0$ and there is nothing to prove.

Otherwise, let $\sigma_0,t_0$ be as in Proposition~\ref{prop:concentric} and $\beta,\gamma$
as in Proposition~\ref{prop:radial}. By Proposition~\ref{prop:betaPF2}, $\beta$ is a
nonnegative log-concave sequence on $\Z-t_0$ with interval, finite support. By
Lemma~\ref{lem:tail}, $\gamma$ is nonincreasing with $\{\gamma>0\}$ a down-set and
$\gamma(y)^2\ge\gamma(y-1)\gamma(y+1)$ for all $y$. By Proposition~\ref{prop:radial},
$G(s)=\gamma(\abs{s-\sigma_0})$, and the lattice $\{\abs{s-\sigma_0}:s\in\Z\}$ is contained
in $\Z-t_0$ because $\sigma_0\equiv t_0\pmod1$. Lemma~\ref{lem:radialLC} now gives
$G\in\PFtwo$.
\end{proof}

\begin{corollary}
Corollary~"single-region criterion" of \textup{[c1]}, which covered $79.9\%$ of slices, is
subsumed: condition~(A) at $m=2$ holds for every slice.
\end{corollary}

\section{Why (G) and concavity fail: an identity}\label{sec:why}

The radial picture explains both negative results of this programme at once, and does so by
an identity rather than by a census.

\begin{proposition}\label{prop:whyfail}
With $\beta,\gamma$ as above, for every $y$
\begin{equation}\label{eq:concdeficit}
  2\gamma(y)-\gamma(y-1)-\gamma(y+1)\ =\ \beta(y)-\beta(y-1).
\end{equation}
Hence $G$ is concave at $s$ with $\abs{s-\sigma_0}=y\ge1$ if and only if
$\beta(y)\ge\beta(y-1)$; that is, $G$ fails concavity exactly on the strictly decreasing
part of $\beta$.
\end{proposition}
\begin{proof}
Substitute $\gamma(y-1)=\gamma(y)+\beta(y-1)$ and $\gamma(y+1)=\gamma(y)-\beta(y)$.
\end{proof}

So concavity of a slice sum holds iff its layer profile $\beta$ is nondecreasing on its
support, which it is not in general: $\beta$ is $c_+$ for a concave $c$, and a concave
sequence decreases eventually unless it is constant. This is the mechanism behind
[c1, §"Concavity is false"]. Concretely, for the witness $k(\cdot,2)=(1,3,6,7,6,3,1)$ at
$n=8$, $\mu=(0,3)$, $\lambda=(4,7)$ one has $\gamma=(7,6,3,1)$ and
$\beta=(1,3,2,1)$, so $\beta(2)-\beta(1)=-1$ and concavity fails at $y=2$, i.e.\ at the two
points $\sigma_0\pm2$ --- exactly the failure $2\cdot3<1+6$ recorded there.

Combining Proposition~\ref{prop:whyfail} with Lemma~\ref{lem:flat}:

\begin{corollary}[(G) is dead, with no residual scope]\label{cor:dead}
Suppose two regions $\mathrm R\ne\mathrm R'$ are nonempty, $G_{\mathrm R'}$ is constant and
positive on an interval $J$, and some $s$ with $s-1,s,s+1\in J$ has
$\abs{s-\sigma_0^{(\mathrm R)}}=y\ge1$ where $\beta^{(\mathrm R)}(y)<\beta^{(\mathrm R)}(y-1)$.
Then \textup{(G)} fails for that pair. Region $\mathrm{III}$ contributes a constant
$G_{\mathrm{III}}$ whenever $T_{\mathrm{III}}$ is a single $t$ with $n_t=1$ or $m_t=1$,
and region $\mathrm I$ likewise; so the configuration is not exceptional.
\end{corollary}

This is the kind of reason the registry asks for: it has no quantifier left to shrink. (G)
is not "unverified outside a range" and not "true with a side condition"; it is false for a
structural reason that applies to any shape whose layer profile is non-constant beside a
flat region, and the region decomposition is simply the wrong cut. The right cut is by
\emph{layer} $r$, not by region $t$: the regions slice the sum along $t$, where the family
is a genuine antichain (registry node \texttt{lemma-T-layer-chain}:
not a chain in $29.9\%$ of instances), while the layers slice it along $r$, where the
family is a nested chain of concentric intervals by Proposition~\ref{prop:concentric}.

\begin{remark}[relation to Lemma T]
This does \emph{not} prove Lemma T of [c1, Conj.~2.6] in its stated generality
(arbitrary concave $1$-Lipschitz $\Phi$ and convex $1$-Lipschitz $\Psi$). Concentricity,
Proposition~\ref{prop:concentric}, uses the specific alignment $A=D-n+1$, $C=B+1$ of the
cylindric data; for general $\Phi,\Psi$ the trapezoids need not share a centre. Lemma T
remains open. What is now proved is the statement Lemma T was introduced to imply.
\end{remark}

\section{Verification}\label{sec:verif}

All code is in \verb|scratch/1001-c2/|. Instruments were required to reproduce a known value
or a known failure before any of their output was believed.

\paragraph{Instrument calibration.}
\begin{itemize}
\item The region engine reproduces the known failure of $(*)$ between individual slice
members at $n=6$, $\mu=(0,3)$, $\lambda=(4,5)$, slice $S_\nu=5$ ([c1, §2.1 item 1]):
trapezoids $(1,1,1,1,1)$ on $[5,9]$ for $t=0$ and $\delta_7$ for $t=2$, failures at
$s=6,8$. It also reproduces the individual $\PFtwo$-ness of the four region sums.
\item \textbf{Estimand check.} An \emph{independent} instrument enumerates chains
$\mu\prec\nu\prec\kappa\prec\lambda$ directly using the vendored primitives
\verb|is_shape|, \verb|hstrip|, \verb|size| of \verb|proofs/code-q254/cyl.py| (the code
behind the 2026-09-20 M-convexity paper), never touching $a_t,b_t,c_t,d_t$. Over
$2\le n\le7$, $d\le11$: the symmetry control $k(a,b)=k(b,a)$ holds on $4908$ of $4908$
entries, and $G(u+a)$ agrees with $k(a,b)$ on all $3427$ (shape, $b$) pairs. So the object
being proved log-concave is $k(\cdot,b)$ and not a near-miss of it.
\item \textbf{Refusal control for Lemma~\ref{lem:tail}.} A first control --- $\beta$
log-concave but not truncated-concave --- produced $0$ failures in $5578$ trials, i.e.\ it
did not fire. That was a signal that the lemma was stronger than the concavity-based proof
I first wrote, not that the instrument was fine; it led to the geometric-decay proof above.
The corrected control --- $\beta$ with interval support but \emph{not} log-concave ---
refuses $309\,362$ of $374\,527$ trials ($82.6\%$), and small explicit probes
($\beta=(10,1,10)$, $(2,1,3)$, $(1,0,5)$) all fire. The instrument can say no.
\end{itemize}

\paragraph{Results.}
\begin{itemize}
\item \textbf{Crossed-sum identities} (Lemma~\ref{lem:crossed}) and concentricity
(Proposition~\ref{prop:concentric}), together with concavity and $1$-Lipschitzness of $L$
(Lemma~\ref{lem:tent}): $0$ violations over $292\,290$ slices, $2\le n\le12$, $d\le14$.
\item \textbf{Radial form} (Proposition~\ref{prop:radial}): $G=\gamma(\abs{\cdot-\sigma_0})$
on all $27\,522$ nonzero slices of the reduced sweep ($0$ mismatches). An earlier run
reported $9616$ mismatches; the defect was in my own tail accumulator, which built $\gamma$
only down to $\min\operatorname{supp}\beta$ instead of to the bottom of the lattice, so
$\gamma(0)$ was never computed when $\beta(0)=0$. Smallest exposing instance $n=4$,
$\mu=(0,1)$, $\lambda=(0,3)$, $b=0$, where $G=(1,1,1)$ and the buggy reconstruction dropped
the middle entry.
\item \textbf{$\beta$ is $\PFtwo$} (Proposition~\ref{prop:betaPF2}): $0$ failures on
$27\,522$ instances; the explicit concave $c$ of~\eqref{eq:c} reproduces $\beta$ exactly on
all $27\,522$ and is concave on the lattice in all $27\,522$. Control: deleting the clamp
$N(\Lambda+1-r)$ makes the formula wrong in $7245$ of $13\,761$ instances, smallest $n=4$,
$\mu=(0,1)$, $\lambda=(0,2)$, $b=1$ --- so the clamp is load-bearing. The clamp is
\emph{active} (ghost rows above $\Lambda$) in $14\,490$ of $27\,522$ instances.
\item \textbf{Lemma~\ref{lem:tail}}: $0$ failures on $161\,771$ random $\PFtwo$ sequences
and on an exhaustive enumeration of all $9654$ $\PFtwo$ sequences in $\{0,\dots,7\}^{L}$,
$L\le6$; plus geometric and long flat sequences up to length $30$.
\item \textbf{Condition (A)}: $0$ failures over $201\,147$ nonempty slices,
$2\le n\le11$, $d\le14$. Extended afterwards, by the direct trapezoid sum
\emph{without} using the radial shortcut, to $2\le n\le 16$, $d\le20$
($89\,366$ nonzero slices, $\mu_1=0$ by translation invariance): $0$ failures. The
$\PFtwo$ test itself refuses $(1,1,2,1,1)$, so it can say no.
\item \textbf{(G)}: $92$ failures over the same range; see Remark~\ref{rem:census}.
\end{itemize}

\section{Gaps, and what $m\ge3$ needs}

\begin{enumerate}
\item \textbf{No gap in Theorem~\ref{thm:main}.} Every step is Lemma~\ref{lem:crossed},
Proposition~\ref{prop:concentric}, Lemma~\ref{lem:tent},
Proposition~\ref{prop:radial}, Proposition~\ref{prop:betaPF2} (which cites
[c1, Lem.~2.2], \texttt{lean-verified}), Lemma~\ref{lem:tail}, Lemma~\ref{lem:radialLC}.
The one external input is [c1, Prop.~2.3], recorded \texttt{proved} in the registry.
\item \textbf{Lemma T is still open} and is now known to be strictly stronger than what (A)
at $m=2$ needs; see the remark at the end of §\ref{sec:why}.
\item \textbf{$m\ge3$.} The proof used $m=2$ twice, both times in
Lemma~\ref{lem:crossed}: there are exactly two beads, so there are exactly two crossed sums,
and each collapses. At general $m$ the slice sum is
$\sum_\nu\mathop{\text{\Large$*$}}_{i}\mathbf 1_{[L_i(\nu)-\nu_i,R_i(\nu)-\nu_i]}$ over an
$(m-1)$-dimensional slice of a box, so the relevant question is whether the summands are
again concentric. The natural generalisation to test first is whether
$\sum_i\bigl(L_i(\nu)+R_i(\nu)\bigr)$ is independent of $\nu$ on each slice. That is a
one-line computation, and it is the next thing I would check.
\item \textbf{What this says about condition (B) and (L3).} Theorem~\ref{thm:main} settles
the $m=2$ half of the $\ell=3$ coupling gap (registry node \texttt{ell3-coupling-gap});
condition (B) (\texttt{conj-B-logsubmodular}) is untouched by this argument and remains
\texttt{computed}.
\end{enumerate}

\section{Summary}

The brief asked me to prove (G) and warned me, correctly, to extend the range first. The
range extension refuted (G) in the first computation of the session, and the refutation was
available all along inside the paper that posed it: against a flat partner the pairwise
inequality \emph{is} concavity, and that paper proves concavity false. The repair was not a
weaker version of (G) but a different decomposition. Cutting the slice sum by layer instead
of by region turns an antichain into a nested chain of concentric intervals, and then
log-concavity is the classical statement that the tail of a log-concave sequence is
log-concave. Condition~(A) at $m=2$ is proved.

\end{document}
